\section{Conclusiones}
\label{sec:conclusiones}

Los dos filtros adaptativos implementados desde cero en este informe
—LMS en WORK\_ELEVEN y NLMS en WORK\_TWELVE— lograron identificar
ciegamente un sistema FIR desconocido a partir únicamente de pares
entrada/salida, recuperando sus coeficientes reales con un error
prácticamente nulo, y cancelar ruido correlacionado tanto sobre una
señal sintética de banda angosta como sobre una grabación de audio
real de contenido espectral más complejo. Estos resultados confirman
experimentalmente, con implementaciones propias y auditables, que el
mecanismo de descenso de gradiente estocástico subyacente a ambos
algoritmos converge de forma robusta en problemas representativos de
identificación de sistemas y cancelación de ruido, sin depender de una
función de filtrado adaptativo ya empaquetada por una librería externa
cuyo funcionamiento interno permanecería opaco.

El análisis de convergencia de LMS frente al tamaño de paso $\mu$
confirmó de forma cuantitativa la cota teórica de estabilidad
$\mu < 2/\lambda_{\max}$: la convergencia observada fue lenta para
$\mu$ muy por debajo de esa cota, rápida y estable para valores
intermedios, y francamente divergente para un $\mu$ ligeramente por
encima del límite teórico calculado para esta configuración
experimental ($\approx 0.4$). Este resultado es relevante más allá de
su valor demostrativo inmediato: evidencia que la elección del tamaño
de paso en un sistema LMS real no puede hacerse de forma arbitraria o
puramente empírica, sino que debe fundamentarse en una estimación,
aunque sea aproximada, de la potencia de la señal de entrada y del
número de coeficientes del filtro.

La comparación directa entre NLMS y LMS bajo una señal de entrada con
potencia no estacionaria constituye el resultado más significativo de
este informe desde el punto de vista práctico: bajo las mismas
condiciones de entrada y el mismo sistema desconocido, LMS con un
$\mu$ fijo —elegido para ser estable en el segmento de mayor
potencia— exhibió un error de identificación aproximadamente 34 veces
mayor que NLMS en el segmento de menor potencia. Esta diferencia no es
un artefacto de la elección particular de parámetros del experimento,
sino la consecuencia directa y esperada de la normalización por
energía instantánea que define a NLMS, y constituye la razón
estructural —no meramente una preferencia de implementación— por la
cual NLMS es el algoritmo de elección en la inmensa mayoría de
aplicaciones reales de filtrado adaptativo sobre señales no
estacionarias en potencia, como la voz humana o los canales de
comunicación con desvanecimiento.

Como trabajo futuro, resultaría natural extender ambos laboratorios
con una variante de paso variable adicional, como LMS normalizado con
factor de olvido (\textit{leaky NLMS}) para mejorar la robustez frente
a sistemas desconocidos que cambian lentamente en el tiempo, o con
RLS (\textit{Recursive Least Squares}) como punto de comparación de
convergencia mucho más rápida a costa de un mayor costo computacional
por iteración —un compromiso distinto al explorado aquí, pero
igualmente relevante en la práctica del procesamiento adaptativo de
señales.
